\scanpage{27}
1.3. Notons que $\unsure{\mathfrak S|\mathscr S}_3$ opère sur $C_n$ de façon
évidente, et les $C_{n'}\text{---}C_n$ sont des $\mathfrak S_3$
morphismes, donc $\mathfrak S_3$ opère sur $C_\infty$.
L'opération de $\mathfrak S_3$ sur $C_1$ \inserted{$\simeq\mathbb P^1_{\mathbb Q}$} est celle
déduite des opérations de $\mathfrak S_3$ sur $\{0,1,\infty\}$.
\inserted{Enfin, $\mathfrak S_3$ opère sur $\unsure{\pi|\Pi}_n$, de façon \deleted{\unsure{que}} que $\mathfrak S_3$ opère
sur $(C_n,\pi_n\ldots)$ -- donc $\mathfrak S_{\unsure{n|3}}$ opère \uline{sur} $(C_\infty,\pi_\infty)$...}

1.4. Considérons la clôture alg. $\overline{\mathbb Q}$ de $\mathbb Q$ dans $\mathbb C$,
d'où par changement de base des
schémas $\overline C_n,\overline C_n^*,\overline C_\infty,\overline C_\infty^*$ sur $\overline C_1=\mathbb P^1_{\overline{\mathbb Q}}$
et $\overline C_1^*=\mathbb P^1_{\overline{\mathbb Q}}\setminus\{0,1,\infty\}$. Sur ces schémas $C_n$
\deleted{\unsure{le groupe}}
\[
\Gamma=\operatorname{Gal}(\overline{\mathbb Q}/\mathbb Q)\tag*{(7)}
\]
opère de façon évidente, et également les
groupes
\[
\left\{\begin{aligned}
\deleted{\unsure{\Gamma|\overline{\Pi}}}\inserted{\unsure{\Pi|\pi}_n=}\deleted{\unsure{\xi|?}}\pi_n(\overline{\mathbb Q})&=
\unsure{\bar\mu|\mu}_n^3/\bar\mu_n\simeq\bar\mu_n\otimes_{\mathbb Z}\mathbb Z^3/\mathbb Z\\
&\simeq\bar\mu_n\otimes_{\widehat{\mathbb Z}}(\widehat{\mathbb Z}^{3}/\widehat{\mathbb Z})\\
\text{où }\quad &\deleted{\unsure{T=\overleftarrow{\mu}_\infty(\overline{\mathbb Q})=\varprojlim_n\mu_n(\overline{\mathbb Q})}}\inserted{\mu_n=\mu_n(\overline{\mathbb Q})}
\end{aligned}\right.\tag*{(8)}
\]
et $\mathfrak S_3$. On voit tout de suite que
ces opérations s'ajustent de façon à ce
qu'opère sur ces courbes le groupe
\[
\unsure{\mathcal E|\Sigma}_n=(\mathfrak S_3\times\Gamma)\underset{1/2}{\mathbin{.}}\Pi_n\tag*{(\deleted{8}9)}
\]
% archive sheet number: 26
